\par\SourceLeaf{125}
Ich bin demnach der Meinung, die sämmtlichen Veränderungsvorschläge des Herrn Universitätsraths O. für den Augenblick ganz auf sich beruhen zu lassen; es ist in dem uns zu lebhaften Danke verpflichtenden P. M. gezeigt, dass eine nahe Gefahr nicht vorhanden ist. Ja selbst wenn in den nächsten Jahren durch noch neu hinzukommende Witwen anstatt des letzten noch immer erfreulichen Ueberschusses einiges Deficit eintreten sollte, so darf man nicht vergessen, dass ja die gesammelten Ueberschüsse zum Theil die Bestimmung haben, solche durch vorübergehende Conjuncturen entstandenen Fluctuationen zu decken.

Aber eine \textit{gründliche} Untersuchung halte ich, in Uebereinstimmung mit dem Rescript und mit den von Sr. Magnificenz geäusserten Ansichten, allerdings für nothwendig. Selbst bei der heitersten Ansicht, die man von dem Zustande der Gesellschaft haben mag, wird eine solche jedenfalls wenigstens späterhin nothwendig werden müssen, schon aus folgendem Grunde.

Wenn ich, \textit{ehe} eine gründliche auf strengen Calcül gegründete Untersuchung statt gefunden hat, meine \textit{Meinung} aussprechen darf, so glaube ich, dass die jetzige grosse Anzahl der Witwen als anomal betrachtet werden muss. Es ist wahr, dass die Anzahl der theilnehmenden Professoren mit der Anzahl der Witwen in einem gewissen Verhältnisse stehen muss; und dass jetzt die erstere Zahl viel grösser ist als ehedem. Allein der \textit{jetzige} hohe Bestand der Witwen steht damit in gar keinem Zusammenhange. Bleibt die Anzahl der theilnehmenden Professoren fortan immer so gross, so ist dies \textit{ein sehr ernsthaft zu erwägender Punkt}, aber nicht für jetzt, sondern wegen der fernen Zukunft; erst nach 20 oder 30 und mehreren Jahren können die Folgen \textit{davon} sehr sichtbar werden.

Dies vorausgesetzt, ist mit Wahrscheinlichkeit anzunehmen, dass, vielleicht schon nach wenigen Jahren, die Anzahl der Witwen wieder abnehmen, vielleicht bedeutend abnehmen, und also der Bestand der Ueberschüsse von 8149 Thl. wieder anwachsen, vielleicht bedeutend anwachsen wird. \textit{Ob aber die Anzahl der Witwen z. B. binnen \textup{10} Jahren bis auf oder unter \textup{15} abnehmen wird, ist viel ungewisser.} Gesetzt nun die Ueberschüsse wären auf 15000 Thl. oder höher angewachsen, die Anzahl der Witwen aber bliebe hartnäckig auf 16 stehen, was soll dann geschehen? Von der einen Seite will man die Ueberschüsse nicht ins Unendliche anwachsen lassen, von der andern steht das \textit{Statut} einer Vergrösserung der Pension entgegen. Dann \textit{muss} ja eine gründliche Prüfung angestellt werden, ob und \textit{in welchem Maasse} man das \textit{Statut verändern} darf, ohne die Gesellschaft zu gefährden.

Dem Vertrauen womit Se. Magnificenz und einige der Herren Collegen mich beehren, indem sie wünschen, dass ich eine solche Prüfung auf mich nehme, durch welche nemlich eine auf Mortalitätsgesetze und die Wahrscheinlichkeitsrechnung basirte Bilanz zwischen dem Vermögen der Anstalt und ihren Obliegenheiten gezogen werden soll, will ich mich nicht entziehen, muss jedoch folgendes bevorworten.

Erstlich haben von der Langwierigkeit solcher Rechnungen diejenigen Herren eine sehr falsche Vorstellung, welche glauben, dass sie binnen vier Wochen vollendet werden können. Zu einer bestimmten Frist kann ich mich also um so weniger anheischig machen, je kleiner der Theil meiner Zeit sein wird, den ich darauf werde verwenden können.

Zweitens lassen sich die Rechnungen mit Gründlichkeit gar nicht führen, ohne die nöthigen \textit{Data}, wovon zur Zeit \textit{gar Nichts} vorliegt. Worin die erforderlichen Data bestehen, werde ich weiterhin angeben; ohne sie kann ich mich auf gar nichts einlassen; ob, auf welche Weise und wie bald sie aber zusammen zu bringen sind, muss ich ganz der Kirchen-Deputation überlassen.

Drittens, die eigenthümliche Einrichtung unsrer Witwenkasse enthält mehrere Elemente die von dem Mortalitätsgesetze unabhängig sind, und sich einem Calcül nicht unterwerfen lassen. Wegen dieses Umstandes wird das Endresultat nothwendig mit einiger Unvollkommenheit behaftet bleiben; ich hoffe jedoch, dass sich Surrogate finden lassen, durch deren Benutzung diese Unvollkommenheit unerheblich sein wird.

Ich will nun suchen, in der Kürze anzudeuten, worauf es bei dieser Arbeit ankommt.
